\documentclass[a4paper,11pt]{article} % type
\usepackage[french]{babel} % règles typographiques françaises
\usepackage{exptech} % mise en forme technique
\usepackage{graphicx} % insérer des images (\includegraphics)
\usepackage{booktabs} % Améliore la qualité des tableaux (lignes horizontales propres)
\usepackage{array} % Ajoute des options pour la mise en forme des tableaux
\usepackage{hyperref} % Ajoute des liens cliquables
\usepackage{listings} % Permet d'insérer et formater du code source
\usepackage{xcolor} % Gère les couleurs (utile pour texte, tableaux, code, etc.)

\lstset{
    basicstyle=\ttfamily\small,
    breaklines=true,
    frame=single,
    backgroundcolor=\color{gray!10}
}

\title{
\begin{center}
{\Huge \textbf{Rapport}} \\[1cm]
{\textbf{} Aide à la lecture pour les manuels techniques d'appareils électroniques}
\end{center}
}

\author{Victoire, Ilyas, Ma\"iwenn, Rayane, Louis 
  \\ \\
  Encadrant : Jean-Louis Pazat}

\date{2025-2026}

\begin{document}

\maketitle

\section*{Introduction}

Un manuel de service pour un amplificateur ou autre appareil électronique ancien contient 
plusieurs dizaines de pages hétérogènes (et souvent mal numérisées) : textes descriptifs, tableaux de composants, 
schémas électroniques et illustrations de cartes.
Pour un passionné souhaitant localiser le condensateur C14, par exemple, il doit
manuellement parcourir la table, retenir la valeur, puis chercher visuellement dans les
schémas et cartes aux styles variés.

L'objectif de ce projet est d'automatiser cette navigation.
Le système doit, à partir d'un PDF, permettre à l'utilisateur de cliquer sur un composant pour le voir mis en surbrillance
dans toutes les vues simultanément.

\section{Choix de la pipeline}
(1)~extraire automatiquement les images pertinentes,
(2)~reconnaître les identifiants de composants dans chacune, et
(3)~lier le même composant en le mettant en surbrillance dans toutes les vues

\section{État de l'art des outils}

La première phase du projet a été consacrée à une exploration large des solutions
disponibles pour extraire du texte depuis des documents PDF contenant des schémas.
On distingue trois familles d'outils : les services en ligne, les bibliothèques OCR
locales, et les détecteurs d'objets par apprentissage profond.

\subsection{Services en ligne}

Les services en ligne ont été les premiers évalués pour leur accessibilité.
Ils produisent généralement un PDF sélectionnable ou un fichier texte, mais présentent
tous des limitations rédhibitoires pour notre cas d'usage.

\paragraph{PDFCandy et Smallpdf}
Ces plateformes web~\cite{pdfcandy}~\cite{smallpdf} proposent une extraction de texte directe depuis
un PDF.
Les noms de nombreux composants sont reconnus, le texte est rendu sans ordre spatial et le résultat devient inutilisable sur les pages mêlant schémas et texte.
Smallpdf présente les mêmes défauts, avec l'inconvénient supplémentaire
d'être propriétaire, non entraînable, et disponible uniquement en ligne (donc
inutilisable dans un pipeline local).

\paragraph{ILovePDF OCR et Adobe Acrobat}
ILovePDF~\cite{ilovepdf} et Adobe Acrobat produisent un PDF sélectionnable
(score~: 6/10), ce qui facilite en théorie l'extraction de liens.
Cependant, les deux échouent sur le texte incliné, omniprésent dans les schémas
électroniques (étiquettes de composants souvent à 90° ou 45°).
Adobe Acrobat reconnaît bien les noms de composants en position droite,
mais reste sans solution pour les orientations non standard.
De plus, ILovePDF est un service payant pour un usage intensif.

\paragraph{ImagetoText}
Cet outil en ligne (score~: 7/10) est partiellement ouvert via une API.
Il reconnaît un grand nombre de noms de composants, mais le rendu textuel est
désordonné : les coordonnées spatiales ne sont pas conservées, ce qui rend
impossible toute correspondance entre un identifiant et sa position dans l'image.

\paragraph{pdf24}
pdf24~\cite{pdf24} obtient un score de 6,5/10.
Son principal défaut est un dictionnaire trop complet : l'outil cherche à lire
tous les caractères du document, y compris les traits et annotations des schémas,
ce qui génère beaucoup de faux positifs.
Il échoue également sur les caractères inclinés.

\subsection{Bibliothèques OCR locales}

\paragraph{Tesseract}
Tesseract~\cite{tesseract} est un moteur OCR open-source maintenu par Google.
Il excelle sur les images de bonne résolution contenant du texte structuré
Il est possible de restreindre son vocabulaire ou de le fine-tuner sur des polices
spécifiques.
En revanche, il reste sensible aux petites polices et aux orientations non standard.

\paragraph{EasyOCR}
EasyOCR~\cite{easyocr} gère mieux les textes inclinés ou de faible contraste que Tesseract,
mais son comportement est moins prédictible et ses hyperparamètres moins accessibles.
Ses résultats sur les étiquettes de composants sont comparables à Tesseract,
sans avantage décisif.

\subsection{Détection d'objets par apprentissage profond}

\paragraph{YOLO v8}
YOLOv8~\cite{yolov8} (\textit{You Only Look Once}, version~8) est un détecteur d'objets.
Sa principale force pour ce projet est son entraînabilité : en fournissant un dataset
annoté au format Roboflow~\cite{roboflow}, on peut lui apprendre à reconnaître des classes
spécifiques (table, schéma, carte, ou zones de texte composant).
Son API Python est simple et s'intègre facilement dans un pipeline.

\subsection{Tableau comparatif}

\begin{center}
\begin{tabular}{lcccc}
\toprule
Outil & Type & Entraînable & Coord. & Texte incliné \\
\midrule
PDFCandy       & En ligne  & Non     & Non & Non \\
Smallpdf       & En ligne  & Non     & Non & Non \\
ILovePDF OCR   & En ligne  & Non     & Non & Non \\
Adobe Acrobat  & En ligne  & Non     & Non & Non \\
ImagetoText    & En ligne  & Partiel & Non & Partiel \\
pdf24          & En ligne  & Non     & Non & Non \\
Tesseract      & Local     & Partiel & Oui & Non \\
EasyOCR        & Local     & Non     & Oui & Partiel \\
YOLO v8        & Local     & Oui     & Oui & --- \\
\bottomrule
\end{tabular}
\end{center}

Aucun outil pris isolément ne répond au besoin complet
Nous partons sur une approche hybride détaillée à la section suivante.

\section{Démarche et obstacles rencontrés}

\subsection{Phase 1 : exploration (septembre -- octobre 2025)}

Les premières séances ont été consacrées à l'exploration individuelle des outils disponibles.
Une question centrale est rapidement apparue : fallait-il un seul outil polyvalent ou
une combinaison spécialisée~?

YOLO a été testé pour détecter directement les \textit{symboles} de composants
(résistances, condensateurs) dans les schémas.
Les résultats ont été décevants : les symboles électroniques sont trop petits, trop proches
les uns des autres, et trop variables selon les fabricants pour constituer des classes
stables sans un dataset très large.

Tesseract a en revanche montré des résultats satisfaisants sur les tables de composants, qui contiennent du texte bien imprimé et structuré.
Un script de post-traitement Excel permettait de corriger les confusions fréquentes
(\texttt{Ç}→\texttt{C}, \texttt{0}→\texttt{O}, etc.).

\subsection{Phase 2 : convergence vers YOLO + Tesseract (novembre 2025)}

À l'issue de la comparaison, nous avons arrêté l'architecture suivante :
\begin{itemize}
  \item \textbf{YOLO} pour détecter les \textit{zones de texte} (séries de lettres et chiffres)
        dans les schémas et cartes~;
  \item \textbf{Tesseract} pour lire le contenu textuel de chacune de ces zones
        et de la table de composants
\end{itemize}

Un datasaset à été construit : 800 boîtes annotées sur des schémas réels. Les annotations ont été effectuées via Roboflow en suivant le format YOLO
(\texttt{train/}, \texttt{valid/}, \texttt{data.yaml}).

\subsection{Phase 3 : problèmes d'intégration (décembre 2025 -- janvier 2026)}

Nous avons rencontré plusieurs points :

\paragraph{Conversion PDF → images}
La bibliothèque \texttt{pdf2image}~\cite{pdf2image} a été retenue pour convertir les pages du PDF
en images.
Le paramètre \texttt{dpi=300} s'est avéré important : en dessous de 200 DPI, les caractères des étiquettes de composants devenaient illisibles pour l'OCR.

\paragraph{Orientation des images}
Certains texte sont à la vertical, Après l' extraction de YOLO, ces images arrivaient pivotées.
Une correction de rotation systématique a été intégrée.

\subsection{Phase 4 : choix entre YOLO et EasyOCR (janvier -- février 2026)}

EasyOCR a été réévalué comme alternative à Tesseract + YOLO pour la lecture des schémas.
Si ses résultats bruts étaient comparables, il a été écarté au profit de Tesseract + YOLO pour
la détection de zones, car YOLO est entraînable et permet de cibler précisément les
étiquettes de composants, là où EasyOCR ne disposait pas d'équivalent pour ce
niveau de personnalisation.

La décision finale a été de conserver détection de zone par YOLO, puis OCR par Tesseract. Deux modèles YOLO sont entraînés ; un pour détecter les tableaux et les schémas, et l'autre pour détecter les étiquettes de composants dans un shéma

\subsection{Phase 5 : intégration finale (mars -- avril 2026)}

Les différentes parties ont été progressivement intégrées dans un Jupyter Notebook
unique (\texttt{NotebookEP.ipynb}).

Le notebook à été transformé en application web via Voilà~\cite{voila}.

\section{Architecture finale du pipeline}

Le pipeline final est composé de trois modules Python orchestrés dans le notebook.

\subsection{Vue d'ensemble}

\begin{center}
\begin{tabular}{clp{6cm}}
\toprule
Étape & Fichier & Rôle \\
\midrule
1 & \texttt{prediction.py} & Extraction brute de toutes les zones par YOLO \\
2 & \texttt{merge.py}      & Sélection de la meilleure zone par classe, rotation, lancement \\
3 & \texttt{test1.py}      & Interface interactive (OCR + affichage) \\
\bottomrule
\end{tabular}
\end{center}

\subsection{\texttt{prediction.py} --- Extraction des zones}

Ce module charge le modèle YOLO entraîné (\texttt{best.pt}) qui reconnaît trois classes :
\texttt{Table}, \texttt{Diagram} et \texttt{Board}.
Il convertit le PDF page par page (DPI~300)
Chaque zone détectée est découpée dans l'image puis sauvegardée en PNG dans le répertoire de sortie.

Le seuil de confiance volontairement bas (0,212) est justifié par le faible taux de faux
positifs du modèle spécialisé : il vaut mieux ne pas manquer une table
qu'écarter une détection incertaine.

\subsection{\texttt{merge.py} --- Orchestration}

Pour chaque classe \texttt{prediction.py}, il ne retient que \textbf{la zone de plus grande surface},
ce qui correspond généralement à l'image la plus exploitable du document.

Il applique ensuite les corrections de rotation et lance l'interface.
L'orchestration est encapsulée dans une fonction \texttt{process\_and\_visualize(pdf\_path, model\_path)}.

\subsection{\texttt{test1.py} --- Interface interactive}

Il charge les images extraites.

\subsubsection{Reconnaissance dans la table}

La table BOM est traitée par Tesseract seul.
On ne lit que la colonne de gauche, car c'est là que se trouvent systématiquement les identifiants.
Avant l'OCR, l'image est convertie en niveaux de gris, agrandie d'un facteur~2
et binarisée pour maximiser le contraste.

\subsubsection{Reconnaissance dans le schéma et la carte}

Pour les schémas et cartes PCB, les étiquettes sont dispersées et très petites.
Une approche hybride est de nouveau utilisée :
\begin{enumerate}
  \item \textbf{YOLO \texttt{bestr.pt}} localise les zones probables d'étiquettes
        (\textit{bounding boxes} avec \texttt{conf=0.20}, \texttt{imgsz=1280})~;
  \item \textbf{Tesseract} lit chaque zone découpée après agrandissement (\(\times 3\))
\end{enumerate}

\subsubsection{Nettoyage des identifiants OCR}

L'OCR introduit des erreurs récurrentes sur les petites polices.
La fonction \texttt{clean\_ocr\_text()} applique un pipeline en trois étapes :
\begin{enumerate}
  \item Suppression des préfixes fantômes (\texttt{AR4} → \texttt{R4})~;
  \item Suppression des zéros de remplissage (\texttt{Q010} → \texttt{Q10})~;
  \item Correction des confusions visuelles (\texttt{RS} → \texttt{R5}).
\end{enumerate}
Les préfixes valides reconnus sont : R, C, L, Q, T, B, S, U, IC, TR, D, Z.

\subsubsection{Interface Matplotlib}

La figure finale présente trois panneaux côte à côte (table, schéma, carte).
La liste des composants de la table est affichée dans un quatrième panneau
sous forme de texte cliquable, colorié par famille (R, C, Q, etc.).

L'interaction utilisateur repose sur trois événements :
\begin{itemize}
  \item \textbf{Clic gauche sur la liste} : zoom sur le composant dans les trois images
        et affichage d'un rectangle rouge~;
  \item \textbf{Clic gauche sur une image} : sélection du composant le plus proche et propagation aux autres panneaux~;
  \item \textbf{Clic droit} : réinitialisation de toutes les vues.
  \item \textbf{Molette} : zoom sur le curseur.
\end{itemize}

\section*{Bilan}

\subsection*{Points positifs}

Le pipeline fonctionne de bout en bout sur les fichiers de test.
L'approche hybride YOLO + Tesseract permet d'allier la précision de localisation du
réseau de neurones à la robustesse de l'OCR classique sur le texte structuré.
Le taux de rappel du modèle de détection de noms atteint 0,80 à un seuil de confiance
de 0,40, ce qui est satisfaisant pour un prototype.

\subsection*{Pistes d'amélioration}

Plusieurs limites ont été identifiées :
\begin{itemize}
    \item Les dataset d'entraînement restent de taille limité.
        Un enrichissement améliorerait le rappel sur les documents non vus.
    \item La transformation en application par l'intermédiaire de Voilà n'est pas finalisée : de fait, l'intervention de clics de l'utilisateur pose des problèmes de synchronisation. 
        
    \item à continuer
\end{itemize}

\section*{Conclusion}

En conclusion, ce projet nous a avant tout permis de manipuler des modèles d'OCR. Il nous a aussi poussés à explorer la manipulation de chaînes de caractères avec Python ainsi que la travail de mise en commun de différents bouts de code.

\bibliography{BiblioEP}

\end{document}
